\subsection{抽象化とカプセル化}
抽象化は、問題解決に不可欠な技法である。対象のすべての細部を同時に扱うのではなく、目的に関係する情報だけに注目する。抽象化は曖昧にすることではない。むしろ、どの水準で正確に議論するかを決めることである。
\subsubsection{抽象化の水準}
抽象化では、全体像の一つの水準に集中する。たとえば、Webサービスを考えるとき、利用者体験の水準、業務処理の水準、APIの水準、データベースの水準、物理サーバの水準は異なる。今どの水準で議論しているかを明確にしないと、議論が混乱する。
抽象化の水準は、現実の部品と一対一に対応するとは限らない。むしろ、標準化されたインタフェースを境界として設けることが重要である。APIのような標準インタフェースは、移植性、統合の容易さ、再利用性を高める。
\subsubsection{カプセル化}
カプセル化は、抽象化を実現する仕組みである。ある水準で作業するとき、その上下の水準の詳細を隠す。利用者は、対象の内部表現や処理手順を知らなくても、提供された操作を通じて対象を使える。
オブジェクト指向で、オブジェクトの内部データを隠し、公開されたメソッドだけを通じて操作するのは典型例である。カプセル化により、内部実装を変更しても、外部インタフェースが保たれていれば利用側への影響を抑えられる。
\subsubsection{階層}
抽象化の水準は、しばしば階層として整理される。階層は、逐次的な一列の構造、木構造、多対多の構造を取ることがある。ただし、階層には循環があってはならない。循環すると、どの水準がどの水準を支えるかが不明確になる。
タスク分解も階層を作る。組織の大きな目標を小さなタスクへ分け、それぞれをさらにサブタスクへ分けると、作業の責任範囲と依存関係が理解しやすくなる。
\subsubsection{代替的な抽象化}
同じ問題に対して、複数の抽象化を同時に持つことが有用な場合がある。たとえば同じソフトウェアについて、クラス図、状態遷移図、シーケンス図を使うことがある。これらは上下関係ではなく、異なる視点から同じ対象を説明する図である。
代替的な抽象化は理解を助ける一方で、同期が難しい。クラス図で変更した内容が状態遷移図やシーケンス図にも反映されているかを確認しなければならない。
\par\smallskip
\begin{center}
\centering
\IfFileExists{assets/generated_figures/ch18/fig-ch18-04.png}{\includegraphics[width=0.96\linewidth]{assets/generated_figures/ch18/fig-ch18-04.png}}{\fbox{\parbox[c][48mm][c]{0.9\linewidth}{\centering 画像生成待ち\\抽象化の水準と代替的な抽象化の図}}}
\par\smallskip
\refstepcounter{figure}\label{fig:ch18-04}図 \thefigure: 抽象化の水準と代替的な抽象化の図
\end{center}
図 \thefigure は、抽象化の水準と代替的な抽象化の図を視覚的に整理したものである。
\par\smallskip
